% Generated draft: requires mathematical and source review.
\documentclass[11pt]{article}
\usepackage[margin=25mm]{geometry}
\usepackage{fontspec}
\setmainfont{Latin Modern Roman}
\usepackage{amsmath,amssymb,mathtools,unicode-math}
\setmathfont{Latin Modern Math}
\usepackage{enumitem,needspace,xurl,hyperref}
\setlength{\parindent}{0pt}
\setlength{\parskip}{5pt}
\setlength{\emergencystretch}{3em}
\newcommand{\N}{\mathbb{N}}
\newcommand{\Z}{\mathbb{Z}}
\newcommand{\Q}{\mathbb{Q}}
\newcommand{\R}{\mathbb{R}}
\newcommand{\C}{\mathbb{C}}
\newcommand{\F}{\mathbb{F}}
\newcommand{\A}{\mathbb{A}}
\newcommand{\PP}{\mathbb{P}}
\newcommand{\E}{\mathbb{E}}
\newcommand{\eps}{\varepsilon}
\newcommand{\dd}{\,\mathrm{d}}
\begin{document}
\section*{UnsolvedMath ID: 800105}
\textbf{Source problem code:} catalogue entry
\subsection{Definitions and mathematical statement}
The supplied record asks for a characterization of pairs of prelaminations with a universal property. It does not define a prelamination, specify the admissible pairs, construct their corresponding bifoliated plane, or define an ``Ansov-like'' action (the spelling in the supplied question).

To display the intended quantifiers without supplying missing hypotheses, let \((L_1,L_2)\) denote a candidate pair, write \(P(L_1,L_2)\) for its corresponding bifoliated plane when that correspondence is defined, and let \(\mathcal A_G(L_1,L_2)\) denote the group actions of \(G\) that preserve the pair in the intended sense. These symbols are placeholders for constructions and classes not specified in the record. The question is to characterize the pairs for which
\[
\forall\,\text{groups }G\;\forall\alpha\in\mathcal A_G(L_1,L_2),\qquad
\alpha\text{ induces an ``Ansov-like'' action of }G\text{ on }P(L_1,L_2).
\]
Here ``induces'' retains the unspecified correspondence in the source; no construction or uniqueness of the induced action is asserted. The admissible prelaminations, the meaning of preservation, the correspondence with a bifoliated plane, and the action property must be supplied before this schema becomes a precise mathematical statement. In particular, the record does not determine whether ``Ansov-like'' means ``Anosov-like,'' or which definition of that term is intended.
\subsection{Short English statement}
Characterize pairs of prelaminations such that every group action preserving the pair induces an ``Ansov-like'' action on the associated bifoliated plane. The necessary definitions and correspondence are absent from the supplied record, so this is a provisional formulation requiring mathematical review.
\subsection{Sources}
\begin{itemize}
\item Supplied local UnsolvedMath registry record: permanent ID \(800105\), exact source problem code \texttt{catalogue entry}, title ``Foliations and Diffeomorphism Groups --- Question 46.'' The record attributes the problem to Nariman and Yazdi, with proposed year 2025.
\item Nariman and Yazdi, \emph{Problem List on Foliations and Diffeomorphism Groups} (2025), Question 46. Supplied source URL: \url{https://arxiv.org/abs/2503.18274}. This is the authoritative-source reference supplied by the registry; its contents have not been independently consulted for this draft.
\item Supplied provenance: access date 2026-10-16; extraction method \texttt{arxiv\_latex}; upstream revision \texttt{d371ee917ac9647fbb06cd6fc529deec6f3dda82}. The access date records supplied provenance, not access performed for this draft. The registry labels the problem ``open''; that label is supplied metadata, not an independently verified literature assessment.
\end{itemize}
\end{document}
